\section*{Introduction}

The point of this course is to \emph{reformulate Newtonian dynamics}, and in several ways. But the first question is, why?

\begin{itemize}
    \item This will allow us to handle more complicated systems. For example, if we have:
    \begin{itemize}
        \item many forces;
        \item constraints;
        \item many particles; or
        \item constrained bodies.
    \end{itemize}
    \item This will emphasise the \emph{role of symmetry} in physics. The type of allowed interactions between fundamental particles is highly constrained by symmetry, and this is a very important idea in modern theoretical physics.
    \item Using \emph{Hamiltonian mechanics}, for example, will allow us to make very close contact to many different areas of maths and physics. For example, the Hamiltonian also appears in \textit{IB Quantum Mechanics}, and recall that \(\comm{\hat{x}}{\hat{p}} = i \hbar\). But this is, fundamentally, not very quantum! We shall see that the \emph{Poisson bracket} \(\pb{x}{p} = 1\) is very closely related. For example, there is a relation \(\comm{L_i}{L_j} = i \hbar \levicivita_{ijk} L_k\), and there is an analogous relation in classical mechanics. This allows us to \emph{understand the structure of quantum mechanics} much better.
    \item There is close connection to \emph{symplectic geometry}, where we of ten study distances using a \emph{symmetric metric}. For example, the standard Euclidean metric is the Kronecker delta \(\delta_{ij}\), while in polar coordinates, it will be \(\dd{s}^2 = \dd{r}^2 + r^2 \dd{\theta}^2\), and there will be a metric tensor in \emph{general relativity} \(g_{\mu\nu}\). But in symplectic geometry, we study distances using an \emph{antisymmetric metric}, which is given by an antisymmetric matrix \(\omega_{ab} = - \omega_{ba}\).
\end{itemize}